\documentclass[a4paper]{article} 
\usepackage[a4paper, total={6in, 8in}]{geometry}
%%%% choose one of the following 4 lines to set colour and/or  date style, 
%\usepackage{pa} % Day, num Month year, and colour (use lp2)
\usepackage{fancyhdr}
%\usepackage{fullpage}
\fancyhead{}
\fancyhead[RO,RE]{\bf Bruce A. Desmarais}
\fancyhead[LO,LE]{\bf Research Statement}


\begin{document} 
\pagestyle{fancy}

\vspace*{-.8cm}

My research is driven by a unifying focus that is part substance and part method---developing and applying rigorous quantitative methods to accurately model the complex interdependence inherent in social and political processes. The work immediately following my dissertation was primarily methodological, focused on developing, translating and illustrating the core technical innovations in the domain of inferential network analysis. Network modeling is the methodology that is best suited to understanding the interdependence in social and political processes. My post-tenure research is focused in three areas: methods for inferential network analysis, the development of methods for the causal study of diffusion social processes, and applications in the domains of public policymaking and digital politics. My research is multimodal, with focus on networks, conventional rectangular data, text, an increasingly, images. In terms of methods related to artificial intelligence, in my research I develop and apply tools for network analysis, causal inference, and machine learning. 


My research is motivated by an interest in developing and applying statistical and computational methods that precisely and transparently characterize the complex ways in which political actors are interdependent. Network science is a growing, cross-­disciplinary field that offers a paradigmatic alternative to the methodological individualism that has dominated quantitative social science for the last half-­century. Until the recent emergence of network and other systems based empirical methods, researchers were faced with the choice between buying into the methodological individualism of established quantitative methods or abandoning the framework of quantitative social science. The network scientific methods, including the methodological innovations offered in my research, eliminate this problem. They permit social and political processes to be represented in their full interactive character and embedded in flexible, multi-­scale structures, while retaining the benefits of transparency, replicability and generalizability offered by other quantitative approaches to social research.

I have made sustained contributions in the areas of inferential network analysis, digital communication research, subnational politics and policy, and causal inference with networks. My research has been published in a broad spectrum of disciplines, all highly relevant to artificial intelligence. These include political science, public policy, computer science, statistics, physics, sociology, public administration, and neuroscience. My research has been supported by ten external awards, totaling over \$2,000,000 in support. My research on inferential network analysis, represented by a Cambridge University Press book by that title, as well as an ongoing series of articles in applied social sciences, statistics, and computer science, aims to develop and illustrate methods for conducting rigorous statistical analysis with network data. Network data presents opportunities to learn about the complexity and interdependence in systems. In this collaborative agenda I have developed new methods for studying longitudinal networks, large networks, and networks with citation structure; I have produced illustrative applications in the areas of international relations, legislative politics, and judicial politics; and I have produced multiple software packages that implement statistical methods. 

My more recent publications and work in progress continue to address methods for network analysis, but I focus more broadly on diffusion in social systems and the applied context of policymaking systems. My work on diffusion in networks has been two-pronged. First, I am interested in the general problem of confounding of selection and influence in networks (i.e., do actors form ties because they are alike, and/or do actors become similar because they are tied in a network). In this area I have published research on both experimental and observational methods for doing causal inference regarding contagion in networks, with articles being published in top journals in network science and political science. Second, I am interested in the problem of diffusion network inference. That is, observing the behaviors of a group of actors over time (e.g., social media accounts posting content, governments making policies), how can we infer which actors influence which other actors? Conducting network inference allows researchers to understand network dependence in the context of applications in which networks cannot be directly measured. I have developed new algorithms for network inference and applied them in the area of public policy diffusion. One of the consistent findings from this work is that, contrary to previous findings in the literature, governments do not rely primarily on their geographic neighbors in modeling new public policies. Rather, they draw upon a mix of neighbors, geographies with similar size and economic profiles, an geographies with similar ideological compositions. 

My most recent research agenda addresses digital politics---focused primarily on digital engagement by and with subnational policymakers. It is through this research agenda that I have integrated generative artificial intelligence methods into my research. This project developed from my work on policy diffusion. In that agenda the unit of analysis is the geographic locale (e.g., state, country, city). I was interested in looking at policymakers at a more micro level. In this research, which is supported by a three-year NSF grant titled Digitally Accountable Public Representation, my collaborators and I are collecting data on policymakers' online communication on Twitter/X, Facebook, YouTube, and TikTok. We have developed measures of misinformation dissemination by officials, position-taking related to vaccines, and discussion of election integrity. We have also recently published research on the structure of the national network of subnational policymakers online. This work has led me to develop extensive expertise in the use of methods of natural language processing and image analysis, with an emphasis on recently-developed deep learning approaches. 

I take a lab model approach to my research---an approach that is less common in the social sciences, but I have found to be highly effective in the context of computational social science research. My research is characterized by a regular pursuit of external funding to support a team of co-PIs, postdocs, graduate research assistants, and undergraduate students. I place significant emphasis on mentoring students and postdocs in my lab, and encourage them to develop ambitions in data science generally speaking, not just academic research. I have former PhD and postdoc advisees who work as data scientists at Meta, Amazon, and Westat; in addition to those who are working in postdocs and tenure track academic positions. This team-based approach to computational social science has enabled me to pursue ambitious multi-year research agendas that result in a variety of publications, datasets, and software resources. 

In my future research I plan to continue working on the development, application, and integration of methods for the study of social networks and digital interaction. The problems that continue to excite me relate to the accurate characterization of complex social systems, and the understanding of harmful and extreme online engagement with and by public officials. In particular, I am interested in how we can draw principled causal inference in the contexts of complex interdependence and multimodal communication data. As my research group and agenda grow, I am finding it increasingly feasible to integrate my research thrusts into singular high-impact projects, with my recent NSF grant on the DAPR project serving as an example. My future work would thrive in the iSchool at Illinois, as all of my research relates to solving problems with data, methods, and computing; and my research is highly interdisciplinary. 


\end{document} 
